\section{Stores --- État global}

L'application utilise des \emph{stores} Svelte (\texttt{writable}) pour partager l'état entre composants sans passer par les props. Deux fichiers définissent tous les stores.

\subsection{dataStore.js --- Données du dictionnaire}

\begin{lstlisting}[language=javascript, title=src/lib/stores/dataStore.js]
import { writable } from 'svelte/store';

export const wordData = writable({});
export const phraseData = writable({});
\end{lstlisting}

\begin{description}
  \item[\texttt{wordData}] Contient l'intégralité du dictionnaire, structuré par catégorie grammaticale. Chargé depuis \texttt{static/data.json} au démarrage via \texttt{+layout.server.js} $\rightarrow$ \texttt{+layout.svelte}. Voir l'annexe~\ref{sec:data-schema} pour le schéma complet.

  \item[\texttt{phraseData}] Contient toutes les phrases d'exemple avec traductions et références. Chargé depuis \texttt{static/phrases.json}.
\end{description}

\subsection{uiStore.js --- État de l'interface}

\begin{lstlisting}[language=javascript, title=src/lib/stores/uiStore.js]
import { writable } from 'svelte/store';

export const pinnedElement = writable(null);
export const accentEnabled = writable(true);
export const selectedWord = writable(null);
export const selectedCategory = writable(null);
export const grammarTableData = writable(null);
\end{lstlisting}

\begin{description}
  \item[\texttt{selectedWord}] Le mot actuellement sélectionné dans la sidebar (ex: \texttt{"будинок"}). Mis à jour par \texttt{WordList} au clic.

  \item[\texttt{selectedCategory}] La catégorie du mot sélectionné (ex: \texttt{"nom"}). Toujours mis à jour en même temps que \texttt{selectedWord}.

  \item[\texttt{accentEnabled}] Booléen indiquant si les accents toniques sont affichés. Défaut : \texttt{true}. Basculé par \texttt{AccentCheckbox}.

  \item[\texttt{grammarTableData}] Objet \texttt{\{word, category, infos, wordData\}} transmis à la sidebar de grammaire lors du survol ou du clic sur un span \texttt{.ukr}. Mis à \texttt{null} quand aucun mot n'est survolé (et qu'aucun élément n'est épinglé).

  \item[\texttt{pinnedElement}] Référence DOM de l'élément épinglé (ou \texttt{null}). Quand un élément est épinglé :
  \begin{itemize}
    \item La sidebar de grammaire reste affichée même sans survol.
    \item Les survols d'autres mots ne modifient pas \texttt{grammarTableData}.
    \item Un second clic sur le même élément le désépingle.
  \end{itemize}
\end{description}

\subsection{Diagramme du flux réactif}

\begin{lstlisting}[title=Flux des stores dans l'application]
Utilisateur clique sur un mot (WordList)
  --> selectedWord.set("mot")
  --> selectedCategory.set("nom")
  --> WordDetails reagit ($derived sur les stores)
      --> NounDetails / VerbDetails / ... s'affiche

Utilisateur survole un span .ukr (HtmlContent)
  --> grammarTableData.set({word, category, infos, wordData})
  --> GrammarSidebar reagit ($derived sur grammarTableData)
      --> GrammarTable genere le tableau

Utilisateur clique sur un span .ukr (HtmlContent)
  --> pinnedElement.set(element)
  --> GrammarSidebar ajoute la classe .pinned

Utilisateur coche/decoche "Accents" (AccentCheckbox)
  --> accentEnabled.set(true/false)
  --> HtmlContent reagit ($effect sur $accentEnabled)
      --> applyAccents() re-parcourt tous les .ukr
\end{lstlisting}
